Écris une fonction \texttt{est\_present(nombres, cible)} qui prend une liste d'entiers \texttt{nombres} triés par ordre croissant et une valeur \texttt{cible}, et renvoie \texttt{True} si \texttt{cible} est une valeur de \texttt{nombres}, \texttt{False} sinon.

L'algorithme de recherche dichotomique est obligatoire : on ne doit pas lire plus de la moitié des valeurs de \texttt{nombres}, ni utiliser le test \texttt{cible in nombres}.

Exemples : \texttt{est\_present([1, 2, 3, 4], 2)}, \texttt{est\_present([1, 2, 3, 4], 1)}, \texttt{est\_present([1, 2, 3, 4], 4)} et \texttt{est\_present([1], 1)} valent \texttt{True} ; \texttt{est\_present([1, 2, 3, 4], 5)}, \texttt{est\_present([1, 2, 3, 4], 0)}, \texttt{est\_present([1], 0)} et \texttt{est\_present([], 1)} valent \texttt{False}.
